\documentclass[journal]{IEEEtran}
\usepackage{cite}
\usepackage{amsmath,amssymb,amsfonts}
\usepackage{amsthm}
\usepackage{bm}
\usepackage{graphicx}
\usepackage{algorithm}
\usepackage{algorithmic}
\usepackage{array}

\theoremstyle{plain}
\newtheorem{theorem}{Theorem}
\newtheorem{proposition}{Proposition}
\newtheorem{corollary}{Corollary}
\theoremstyle{definition}
\newtheorem{assumption}{Assumption}

\begin{document}

\title{Measurement-Decoupled Attention--Koopman--T--S Fuzzy Monitoring via Structured Basis-Perturbation Residual Learning}

\author{Anonymous Authors}

\maketitle

\begin{abstract}
This paper addresses normal-data-only monitoring of closed-loop nonlinear systems when physical fault channels, directions, temporal profiles, and labeled fault samples are unavailable. First, a strictly causal attention encoder is coupled with a Takagi--Sugeno (T--S) fuzzy Koopman model to construct a measurement-decoupled normal reference. A start-indexed residual then gives the same mathematical object for free-rollout training, short-horizon monitoring, and long-reference monitoring. Second, the complete fuzzy Koopman Jacobian is decomposed into the interpolated local dynamics and the variation of normalized rule weights. This yields an implementable infinity-norm propagation bound and a reliable reference horizon without eigenvalue, singular-value, inverse-matrix, or higher-order derivative operations during training. Third, model-internal basis perturbations are applied separately to input, measurement, and lifted coordinates. Ordinary forward rollouts produce finite-horizon residual responses used to jointly reduce nominal propagation, increase perturbation sensitivity, and separate response patterns. Online scores are standardized by fuzzy operating region, controller command, command variation, and reference age, and their thresholds are calibrated using held-out normal blocks. Sufficient residual-space conditions are derived for detection and unique group separation. When the conditions for a unique decision fail, the method returns a set of compatible perturbation groups rather than an unsupported physical fault label.
\end{abstract}

\begin{IEEEkeywords}
Attention, fault detection, Koopman operator, normal-data-only learning, structured perturbation, Takagi--Sugeno fuzzy systems.
\end{IEEEkeywords}

\section{Introduction}
\label{sec:intro}

\IEEEPARstart{N}{onlinear} closed-loop systems frequently operate across multiple regimes while providing few representative fault records. Learning normal input--output behavior is therefore attractive for monitoring, but it creates a structural difficulty: the same measurement may be used both to form a residual and to update the state and model against which that residual is evaluated. When an abnormal measurement updates an encoder, an attention context, or a fuzzy scheduling variable, the reference may move with the abnormal trajectory and reduce the discrepancy intended for detection.

Koopman learning seeks lifted coordinates in which nonlinear dynamics admit a finite-dimensional approximation~\cite{brunton2016,lusch2018}. T--S fuzzy models interpolate local linear descriptions through normalized firing strengths~\cite{takagi1985,tanaka2001,lam2018}. Their combination has already been investigated in modeling and predictive control~\cite{hu2026}; the combination itself is consequently not claimed as a contribution here. The T--S structure is instead made operational in three places: the normal reference, the finite-horizon error propagation, and the operating-condition-dependent monitoring rule. Attention is retained to represent long input--output histories, but its output is prevented from feeding new measurements into the normal reference after initialization.

Classical fuzzy fault-detection methods commonly introduce explicit fault or unknown-input channels and then design observers or filters by Lyapunov and linear-matrix-inequality techniques~\cite{nguang2007,yang2011,thumati2014,zhao2020,ran2022}. Koopman-based fault diagnosis has also used learned residual and parity relations~\cite{bakhtiaridoust2023,irani2024}. Those formulations are appropriate when a physical fault model or representative fault data are available. In the setting considered here, training data are normal, and the physical fault location, direction, and time profile are unknown. A physical minimum fault gain cannot therefore be identified. The available structure is limited to the known input and measurement coordinates and to the coordinates of the learned lifted model.

% PLACEHOLDER - Introduction continuation

This information boundary creates one connected design problem. Removing new measurements from the reference prevents adaptation to abnormal observations, but also exposes normal free-rollout drift. Reducing only that drift is insufficient because a less sensitive model may reduce abnormal residuals as well. Finally, an unknown physical fault direction precludes a justified low-dimensional fault-sensitive projection. The method must instead control normal propagation, retain the joint residual, and learn separable responses only for explicitly defined model-coordinate perturbation groups.

The principal contributions are as follows.
\begin{enumerate}
\item A measurement-decoupled attention--Koopman--T--S reference and a start-indexed residual are constructed. They unify free-rollout training, short-horizon monitoring, and long-reference monitoring, and yield an exact nonlinear residual recursion together with a delay-confirmed anchor rule.
\item The complete T--S fuzzy Koopman Jacobian is decomposed to expose the effects of local dynamics, rule overlap, local-model disparity, and scheduling sensitivity. An infinity-norm bound, enforced by elementwise parameterization, gives a finite-horizon normal residual bound and a reliable reference horizon without costly or poorly differentiable matrix operations.
\item Structured input-, measurement-, and lifted-coordinate perturbation rollouts are integrated with nominal propagation learning. The resulting input-conditioned statistic admits sufficient residual-space detection and group-separation conditions and returns a set-valued decision when a unique group is not supported.
\end{enumerate}

The remainder of the paper is organized as follows. Section~\ref{sec:prelim} specifies the information boundary and formulates three aligned subproblems. Section~\ref{sec:method} develops the model, reference, propagation analysis, and structured response learning. Section~\ref{sec:detection} gives the input-conditioned detection and set-valued isolation method. Section~\ref{sec:conclusion} concludes the paper.

\section{Preliminaries and Problem Formulation}
\label{sec:prelim}

\section{Preliminaries and Problem Formulation}
\label{sec:prelim}

\subsection{Information Boundary and Plant Description}

Let $\bm{u}_k \in \mathbb{R}^{m_u}$ be the controller command available to the monitor and let $\bm{u}_k^a$ be the action applied to the plant. The closed-loop plant is described by
\begin{equation}
\begin{aligned}
\bm{u}_k^a &= \bm{u}_k + \bm{f}_k^a, \\
\bm{x}_{k+1} &= \mathcal{X}_0(\bm{x}_k, \bm{u}_k^a, \bm{\xi}_k)\bm{f}_k^p + \bm{\varepsilon}_k^x, \\
\bm{y}_k &= \mathcal{Y}_0(\bm{x}_k, \bm{u}_k^a, \bm{\xi}_k)\bm{f}_k^s + \bm{\varepsilon}_k^y,
\end{aligned}
\label{eq:plant}
\end{equation}
where $\bm{x}_k \in \mathbb{X} \subseteq \mathbb{R}^{m_x}$, $\bm{y}_k \in \mathbb{Y} \subseteq \mathbb{R}^{m_y}$, and $\bm{\xi}_k \in \mathbb{R}^{m_\xi}$ denote the physical state, measured output, and measured exogenous variable, respectively. The healthy mappings $\mathcal{X}_0$ and $\mathcal{Y}_0$ are unknown. The fault vectors in~\eqref{eq:plant} define only possible physical locations; their dimensions do not provide known fault directions to the learning algorithm.

The monitor receives $\bm{u}_k$, $\bm{y}_k$, and $\bm{\xi}_k$. The action $\bm{u}_k^a$ may be unavailable. Using the controller command in the normal model answers the conditional question: how should a healthy plant evolve under the recorded command? A command modified by a feedback controller after an abnormal measurement remains a known conditioning variable, whereas an unmeasured difference between $\bm{u}_k^a$ and $\bm{u}_k$ remains a plant-side inconsistency.

Only normal trajectories are used for model training and threshold calibration. No physical fault matrix, fault direction, fault derivative, fault time model, or labeled fault trajectory is used. Model parameters and calibration quantities are frozen during testing. Define the strictly past history
\begin{equation}
\mathcal{H}_k^- = \{\bm{u}_{k-\ell_h:k-1}, \bm{y}_{k-\ell_h:k-1}, \bm{\xi}_{k-\ell_h:k}\},
\label{eq:history}
\end{equation}
where the current output $\bm{y}_k$ is excluded.

% PLACEHOLDER - Section II-B

\subsection{Finite-History Lifted Representation}

For the healthy dynamics, Koopman evolution acts on an observable by composition with the state transition~\cite{brunton2016}. A finite vector of observables $\bm{\psi}: \mathbb{X} \to \mathbb{Z} \subseteq \mathbb{R}^{m_z}$ would satisfy the ideal composition relation
\begin{equation}
\bm{\psi}(\mathcal{X}_0(\bm{x}_k, \bm{u}_k, \bm{\xi}_k)) = \bm{\psi}(\bm{x}_{k+1}).
\label{eq:koopman_ideal}
\end{equation}

The paper does not assume that a finite invariant subspace exists or that the learned lifted coordinates contain the physical state. Instead, a finite-history encoder estimates a predictive coordinate $\bm{z}_k^{\mathrm{enc}} \in \mathbb{Z}$ from~\eqref{eq:history}, and a T--S interpolation of local lifted linear models approximates~\eqref{eq:koopman_ideal} on the normal operating domain. Approximation and finite-history errors are retained in the residual propagation derived in Section~\ref{sec:method}.

\subsection{Problem Formulation}

Under the preceding information boundary, this paper addresses three subproblems.
\begin{enumerate}
\item \emph{Reference construction:} construct a normal reference that uses new measurements for residual evaluation but not for reference updates, while using one residual definition in training and in both short- and long-horizon testing.
\item \emph{Normal propagation:} quantify the finite-horizon drift of that reference through the complete attention--Koopman--T--S mapping using operations suitable for first-order network training, and make the realized input and fuzzy operating region explicit.
\item \emph{Residual-space decision:} using normal data only, jointly reduce normal residual propagation and increase separability of reproducible model-coordinate perturbation responses; then construct an input-conditioned detection statistic and return a unique group only when the residual response supports it.
\end{enumerate}

\section{Proposed Method}
\label{sec:method}

\section{Proposed Method}
\label{sec:method}

\subsection{Attention--Koopman--T--S Normal Model}

Each element of~\eqref{eq:history} is embedded as a token $\bm{h}_j$. For attention head $q \in \{1, \ldots, m_{\mathrm{hd}}\}$, the strictly causal coefficient assigned at sample $k$ to a past token $j$ is
\begin{equation}
\varpi_{k,j}^{(q)} = \frac{\exp\left((\bm{W}_Q^{(q)} \bm{h}_{k-1})^\mathrm{T} (\bm{W}_K^{(q)} \bm{h}_j) / \sqrt{m_a}\right)}{\sum_{\ell=k-\ell_h}^{k-1} \exp\left((\bm{W}_Q^{(q)} \bm{h}_{k-1})^\mathrm{T} (\bm{W}_K^{(q)} \bm{h}_\ell) / \sqrt{m_a}\right)}.
\label{eq:attention}
\end{equation}

The head outputs are concatenated to form
\begin{equation}
\bm{\chi}_k = \bm{W}_O \left[ \sum_{j=k-\ell_h}^{k-1} \varpi_{k,j}^{(q)} \bm{W}_V^{(q)} \bm{h}_j \right]_{q=1}^{m_{\mathrm{hd}}},
\label{eq:context}
\end{equation}
\begin{equation}
\bm{z}_k^{\mathrm{enc}} = \mathcal{H}_\Theta(\mathcal{H}_k^-).
\label{eq:encoder}
\end{equation}
The encoder in~\eqref{eq:encoder} includes the token embeddings and the multihead attention operation. Equation~\eqref{eq:attention} makes the causal information path explicit; attention itself is not treated as a separate contribution.

The scheduling vector is
\begin{equation}
\bm{\rho}_k = \bm{W}_z \bm{z}_k + \bm{W}_\chi \bm{\chi}_k + \bm{W}_u \bm{u}_k + \bm{W}_\xi \bm{\xi}_k.
\label{eq:scheduling}
\end{equation}
Constant offsets are represented by an augmented constant coordinate. For $i \in \{1, \ldots, m_r\}$, define
\begin{equation}
\begin{aligned}
\varphi_i(\bm{\rho}) &= -\frac{1}{2}(\bm{\rho} - \bm{c}_i)^\mathrm{T} \bm{\Pi}_i (\bm{\rho} - \bm{c}_i), \\
\omega_i(\bm{\rho}) &= \frac{\exp(\varphi_i(\bm{\rho}))}{\sum_{j=1}^{m_r} \exp(\varphi_j(\bm{\rho}))}.
\end{aligned}
\label{eq:membership}
\end{equation}

% PLACEHOLDER - Section III-A continuation

Thus, $\omega_i \geq 0$ and $\sum_i \omega_i = 1$. The diagonal metric is parameterized by
\begin{equation}
\begin{aligned}
\bm{\Pi}_i &= \mathrm{Diag}[\underline{\pi} + (\overline{\pi} - \underline{\pi}) \cdot \mathrm{sigmoid}(\widetilde{\bm{\pi}}_i)], \\
\bm{c}_i &= \overline{c} \cdot \tanh(\widetilde{\bm{c}}_i),
\end{aligned}
\label{eq:param_metric}
\end{equation}
where all operations inside the brackets are elementwise and $0 < \underline{\pi} < \overline{\pi}$. Cross-coordinate dependence is retained through~\eqref{eq:scheduling}; consequently, the diagonal metric in the scheduling space does not imply an axis-aligned partition of the original history space.

The local and interpolated lifted models are
\begin{equation}
\begin{aligned}
\bm{v}_i(\bm{z}, \bm{u}) &= \bm{A}_i \bm{z} + \bm{B}_i \bm{u}, \\
\mathcal{T}_\Theta(\bm{z}, \bm{u}, \bm{\xi}; \bm{\chi}) &= \sum_{i=1}^{m_r} \omega_i(\bm{\rho}) \bm{v}_i(\bm{z}, \bm{u}), \\
\bm{y}^{\mathrm{dec}} &= \mathcal{D}_\Theta(\bm{z}, \bm{u}, \bm{\xi}).
\end{aligned}
\label{eq:ts_model}
\end{equation}
The decoder $\mathcal{D}_\Theta$ is distinct from the physical output mapping $\mathcal{Y}_0$. This distinction is necessary because the learned lifted state need not contain the physical state.

The matrices that determine the propagation bound use elementwise hard parameterizations, for example,
\begin{equation}
\begin{aligned}
[\bm{A}_i]_{pq} &= \frac{\kappa_A}{m_z} \tanh([\widetilde{\bm{A}}_i]_{pq}), \\
[\bm{B}_i]_{pq} &= \frac{\kappa_B}{m_u} \tanh([\widetilde{\bm{B}}_i]_{pq}), \\
[\bm{W}_z]_{pq} &= \frac{\kappa_W}{m_z} \tanh([\widetilde{\bm{W}}_z]_{pq}).
\end{aligned}
\label{eq:hard_param}
\end{equation}
Hence $\|\bm{A}_i\|_\infty \leq \kappa_A$, $\|\bm{B}_i\|_\infty \leq \kappa_B$, and $\|\bm{W}_z\|_\infty \leq \kappa_W$ by construction. No spectral projection is required.

% PLACEHOLDER - Section III-A training

Training uses a normal training split and a free-rollout horizon $\ell_{\mathrm{tr}}$. Starting from $\bm{z}_{s|s}^{\mathrm{fr}} = \bm{z}_s^{\mathrm{enc}}$, the predicted output, rather than the new measurement, updates the rollout history:
\begin{equation}
\begin{aligned}
\bm{\chi}_{t|s}^{\mathrm{fr}} &= \mathcal{H}_\Theta(\mathcal{H}_{t|s}^{\mathrm{fr}}), \\
\bm{z}_{t+1|s}^{\mathrm{fr}} &= \mathcal{T}_\Theta(\bm{z}_{t|s}^{\mathrm{fr}}, \bm{u}_t, \bm{\xi}_t; \bm{\chi}_{t|s}^{\mathrm{fr}}), \\
\bm{y}_{t|s}^{\mathrm{fr}} &= \mathcal{D}_\Theta(\bm{z}_{t|s}^{\mathrm{fr}}, \bm{u}_t, \bm{\xi}_t).
\end{aligned}
\label{eq:rollout}
\end{equation}

The normal prediction loss is
\begin{equation}
\mathcal{L}_{\mathrm{pred}} = \sum_s \sum_{\ell=1}^{\ell_{\mathrm{tr}}} \left( \lambda_z \|\bm{z}_{s+\ell}^{\mathrm{enc}} - \bm{z}_{s+\ell|s}^{\mathrm{fr}}\|_1 + \lambda_y \|\bm{y}_{s+\ell} - \bm{y}_{s+\ell|s}^{\mathrm{fr}}\|_1 \right).
\label{eq:loss_pred}
\end{equation}
Free rollout, rather than one-step teacher forcing, matches the online reference mechanism.

\subsection{Measurement-Decoupled Reference and Unified Residual}

For any start $s$, the rollout~\eqref{eq:rollout} defines the lifted residual
\begin{equation}
\bm{e}_{k|s}^z = \bm{z}_k^{\mathrm{enc}} - \bm{z}_{k|s}^{\mathrm{fr}}, \quad k \geq s,
\label{eq:lifted_residual}
\end{equation}
and the output-consistency residual is
\begin{equation}
\bm{e}_k^y = \bm{y}_k - \mathcal{D}_\Theta(\bm{z}_k^{\mathrm{enc}}, \bm{u}_k, \bm{\xi}_k).
\label{eq:output_residual}
\end{equation}
Equation~\eqref{eq:lifted_residual} is the same object used in~\eqref{eq:loss_pred}. Online, $s = k - \ell_{\mathrm{sh}}$ gives a short rolling residual. Let $s_k^0$ be the latest delay-confirmed normal anchor and define
\begin{equation}
\bm{z}_k^0 = \bm{z}_{k|s_k^0}^{\mathrm{fr}}, \quad \ell_k^0 = k - s_k^0.
\label{eq:long_reference}
\end{equation}

% PLACEHOLDER - Section III-B continuation

Then $s = s_k^0$ gives the long-reference residual. The minimal joint residual is
\begin{equation}
\bm{e}_k = [(\bm{e}_k^y)^\mathrm{T} \; (\bm{e}_{k|k-\ell_{\mathrm{sh}}}^z)^\mathrm{T} \; (\bm{e}_{k|s_k^0}^z)^\mathrm{T}]^\mathrm{T}.
\label{eq:joint_residual}
\end{equation}

Thus, training and testing do not introduce different lifted residuals; only the rollout start changes.

An encoded state is not immediately accepted as a normal anchor. It is stored as a candidate and committed after $\ell_{\mathrm{cf}}$ samples only if no warning is issued for any window containing that candidate. During warning and alarm intervals, the accepted anchor is frozen and~\eqref{eq:rollout} continues without measurement updates.

\begin{proposition}[Delayed anchor acceptance]
\label{prop:anchor}
Suppose every abnormal trajectory in a specified detectable class triggers a warning within at most $\ell_{\mathrm{det}}$ samples after its onset. If $\ell_{\mathrm{cf}} \geq \ell_{\mathrm{det}}$, no candidate generated at or after that onset is committed before the warning.
\end{proposition}

\begin{proof}
A candidate requires $\ell_{\mathrm{cf}}$ consecutive warning-free samples before commitment. Under the stated premise, a warning occurs no later than $\ell_{\mathrm{det}} \leq \ell_{\mathrm{cf}}$ samples after onset, so a post-onset candidate fails the acceptance condition.
\end{proof}

Proposition~\ref{prop:anchor} is conditional; faults that never change the available input--output information cannot be covered. After an alarm clears, the method waits at least $\ell_h$ samples, so that the encoder history no longer contains the flagged interval, and then repeats delay confirmation.

Define the encoded-state innovation and the context mismatch as
\begin{equation}
\begin{aligned}
\bm{\nu}_k &= \bm{z}_{k+1}^{\mathrm{enc}} - \mathcal{T}_\Theta(\bm{z}_k^{\mathrm{enc}}, \bm{u}_k, \bm{\xi}_k; \bm{\chi}_k), \\
\bm{\delta}_{k|s}^\chi &= \mathcal{T}_\Theta(\bm{z}_k^{\mathrm{enc}}, \bm{u}_k, \bm{\xi}_k; \bm{\chi}_k) - \mathcal{T}_\Theta(\bm{z}_k^{\mathrm{enc}}, \bm{u}_k, \bm{\xi}_k; \bm{\chi}_{k|s}^{\mathrm{fr}}).
\end{aligned}
\label{eq:innovation}
\end{equation}

% PLACEHOLDER - Theorem 1

\begin{theorem}[Exact start-indexed residual propagation]
\label{thm:propagation}
Assume $\mathcal{T}_\Theta$ is continuously differentiable in $\bm{z}$ on the line segment joining $\bm{z}_{k|s}^{\mathrm{fr}}$ and $\bm{z}_k^{\mathrm{enc}}$. Then
\begin{equation}
\bm{e}_{k+1|s}^z = \bm{\Gamma}_{k|s} \bm{e}_{k|s}^z + \bm{\nu}_k + \bm{\delta}_{k|s}^\chi,
\label{eq:recursion}
\end{equation}
where
\begin{equation}
\bm{\Gamma}_{k|s} = \int_0^1 \bm{\Gamma}_{z,k}(\bm{z}_{k|s}^{\mathrm{fr}} + \tau \bm{e}_{k|s}^z) \, d\tau.
\label{eq:jacobian_integral}
\end{equation}
Consequently,
\begin{equation}
\bm{e}_{k|s}^z = \bm{\Phi}(k,s) \bm{e}_{s|s}^z + \sum_{j=s}^{k-1} \bm{\Phi}(k, j+1) (\bm{\nu}_j + \bm{\delta}_{j|s}^\chi),
\label{eq:propagation}
\end{equation}
with $\bm{\Phi}(k,j) = \bm{\Gamma}_{k-1|s} \cdots \bm{\Gamma}_{j|s}$ and $\bm{\Phi}(j,j) = \bm{I}$.
\end{theorem}

The proof follows by adding and subtracting the transition evaluated at the encoded state and the rollout context, followed by the integral mean-value identity; details are given in Appendix~A. The term $\bm{\nu}_k$ is the one-step normal-model inconsistency, whereas $\bm{\delta}_{k|s}^\chi$ retains the effect of finite attention memory.

\subsection{Fuzzy-Structure-Dependent Normal Propagation}

For compactness, define $\bm{v}_{i,k} = \bm{A}_i \bm{z}_k + \bm{B}_i \bm{u}_k$ and $\overline{\bm{v}}_k = \sum_i \omega_{i,k} \bm{v}_{i,k}$. From~\eqref{eq:membership},
\begin{equation}
\begin{aligned}
\bm{\gamma}_{i,k}^\rho &:= \nabla_{\bm{\rho}} \varphi_{i,k} = -\bm{\Pi}_i (\bm{\rho}_k - \bm{c}_i), \\
\overline{\bm{\gamma}}_k^\rho &:= \sum_j \omega_{j,k} \bm{\gamma}_{j,k}^\rho, \\
\nabla_{\bm{\rho}} \omega_{i,k} &= \omega_{i,k} (\bm{\gamma}_{i,k}^\rho - \overline{\bm{\gamma}}_k^\rho).
\end{aligned}
\label{eq:gradient}
\end{equation}

% PLACEHOLDER - Theorem 2

\begin{theorem}[Complete fuzzy lifted-state Jacobian]
\label{thm:jacobian}
For fixed $\bm{u}_k$, $\bm{\xi}_k$, and $\bm{\chi}_k$, the Jacobian of~\eqref{eq:ts_model} with respect to $\bm{z}$ is
\begin{equation}
\bm{\Gamma}_{z,k} = \sum_{i=1}^{m_r} \omega_{i,k} \bm{A}_i + \sum_{i=1}^{m_r} \omega_{i,k} (\bm{v}_{i,k} - \overline{\bm{v}}_k) (\bm{\gamma}_{i,k}^\rho - \overline{\bm{\gamma}}_k^\rho)^\mathrm{T} \bm{W}_z.
\label{eq:complete_jacobian}
\end{equation}
Moreover,
\begin{equation}
\|\bm{\Gamma}_{z,k}\|_\infty \leq \kappa_k := \sum_i \omega_{i,k} \|\bm{A}_i\|_\infty + \|\bm{W}_z\|_\infty \sum_i \omega_{i,k} \|\bm{v}_{i,k} - \overline{\bm{v}}_k\|_\infty \|\bm{\gamma}_{i,k}^\rho - \overline{\bm{\gamma}}_k^\rho\|_1.
\label{eq:bound}
\end{equation}
\end{theorem}

Equation~\eqref{eq:complete_jacobian} separates the interpolated local dynamics from the change in normalized firing strengths. The second term becomes large only when rules overlap, their local next-state values differ, and the scheduling vector is sensitive in a direction that changes their relative weights. This is the specific role of the fuzzy structure in the reference analysis. The proof and the induced-norm bound are given in Appendix~B.

\begin{assumption}[Normal residual forcing]
\label{as:forcing}
On a prescribed normal operating domain, known nonnegative values $\overline{\varepsilon}_k$ satisfy
\begin{equation}
\|\bm{\nu}_k + \bm{\delta}_{k|s}^\chi\|_\infty \leq \overline{\varepsilon}_k
\label{eq:forcing_bound}
\end{equation}
on the calibrated normal event.
\end{assumption}

The values in Assumption~\ref{as:forcing} are not claimed as deterministic global bounds obtained from finite data. The structural part of the bound comes from~\eqref{eq:hard_param}--\eqref{eq:bound}; the unresolved normal remainder is calibrated on held-out normal blocks in Section~\ref{sec:detection}.

\begin{corollary}[Finite-horizon normal residual bound]
\label{cor:bound}
If $\beta_{s|s} \geq \|\bm{e}_{s|s}^z\|_\infty$ and
\begin{equation}
\beta_{k+1|s} = \kappa_k \beta_{k|s} + \overline{\varepsilon}_k,
\label{eq:beta_recursion}
\end{equation}
then $\|\bm{e}_{k|s}^z\|_\infty \leq \beta_{k|s}$ while the trajectory remains in the prescribed normal domain. For an allowed reference tolerance $\overline{\beta}$, a reliable horizon is
\begin{equation}
\ell_{\mathrm{ref}}(s) = \max\{\ell : \beta_{s+j|s} \leq \overline{\beta}, \; 1 \leq j \leq \ell\}.
\label{eq:reliable_horizon}
\end{equation}
\end{corollary}

The empirical training counterpart of~\eqref{eq:beta_recursion},
\begin{equation}
\mathcal{L}_{\mathrm{nr}} = \sum_s \sum_{\ell=1}^{\ell_{\mathrm{tr}}} \beta_{s+\ell|s}^{\mathrm{emp}},
\label{eq:loss_nr}
\end{equation}
is used only as a normal-reference propagation regularizer. It is not by itself a claim of improved physical fault sensitivity.

\subsection{Structured Basis-Perturbation Residual Learning}

Reducing~\eqref{eq:loss_nr} controls only the normal side of residual separation. To create a diagnostic training signal without physical fault data, define the model-coordinate group set
\begin{equation}
\mathbb{G}_b = \{u, y, z\},
\label{eq:groups}
\end{equation}
representing input-, measurement-, and lifted-coordinate perturbations. These groups specify where a perturbation is inserted in the learned architecture; they are not physical actuator, sensor, or process labels.

At a sampled start $s$, select one standard basis vector and signed amplitude $\alpha$. The input branch replaces $\bm{u}_s$ by $\bm{u}_s + \alpha \bm{q}_j^u$ in the model rollout. The measurement branch replaces $\bm{y}_s$ by $\bm{y}_s + \alpha \bm{q}_j^y$ in the current output residual and in subsequent causal histories. The lifted branch uses
\begin{equation}
\bm{z}_{s+1}^{(z,j,\alpha)} = \mathcal{T}_\Theta(\bm{z}_s^{\mathrm{enc}}, \bm{u}_s, \bm{\xi}_s; \bm{\chi}_s) + \alpha \bm{q}_j^z.
\label{eq:lifted_pert}
\end{equation}
Thereafter, every branch follows the same frozen model and reference logic as the nominal branch.

For horizon $\ell_b$, stack the branch residuals and subtract the nominal residual stack:
\begin{equation}
\Delta\bm{e}_{s,\ell_b}^{g,j,\alpha} = [(\bm{e}_s^{g,j,\alpha} - \bm{e}_s^0)^\mathrm{T} \; \cdots \; (\bm{e}_{s+\ell_b-1}^{g,j,\alpha} - \bm{e}_{s+\ell_b-1}^0)^\mathrm{T}]^\mathrm{T}.
\label{eq:delta_residual}
\end{equation}

% PLACEHOLDER - Section III-D continuation

Define the scaled response and its unit-$\ell_1$ pattern by
\begin{equation}
\eta_g^{\mathrm{rsp}} = \frac{\|\Delta\bm{e}_{s,\ell_b}^{g,j,\alpha}\|_1}{|\alpha| + \varepsilon_0}, \quad \bm{p}_{g,j,\alpha} = \frac{\Delta\bm{e}_{s,\ell_b}^{g,j,\alpha}}{\|\Delta\bm{e}_{s,\ell_b}^{g,j,\alpha}\|_1 + \varepsilon_0},
\label{eq:response}
\end{equation}
where the absolute value in the numerator is elementwise and $\varepsilon_0 > 0$ is numerical regularization.

The sensitivity and intergroup separation losses are
\begin{equation}
\begin{aligned}
\mathcal{L}_{\mathrm{sen}} &= \sum_{g,j,\alpha} [\overline{\eta}_g - \eta_g^{\mathrm{rsp}}]_+^2, \\
\mathcal{L}_{\mathrm{iso}} &= \sum_{g \neq \overline{g}} [\eta_{\mathrm{iso}} - \|\bm{p}_{g,j,\alpha} - \bm{p}_{\overline{g},\overline{j},\overline{\alpha}}\|_1]_+^2.
\end{aligned}
\label{eq:loss_sens_iso}
\end{equation}
The complete objective is
\begin{equation}
\mathcal{L} = \mathcal{L}_{\mathrm{pred}} + \lambda_{\mathrm{nr}} \mathcal{L}_{\mathrm{nr}} + \lambda_{\mathrm{sen}} \mathcal{L}_{\mathrm{sen}} + \lambda_{\mathrm{iso}} \mathcal{L}_{\mathrm{iso}} + \lambda_{\mathrm{bal}} \mathcal{L}_{\mathrm{bal}},
\label{eq:loss_total}
\end{equation}
where $\mathcal{L}_{\mathrm{bal}}$ penalizes only rule occupancies below a prescribed minimum; it does not force uniform rule usage. Each mini-batch samples one coordinate from each group, so the additional computation consists of three ordinary forward rollouts and first-order backpropagation. No Hessian, eigendecomposition, singular-value decomposition, matrix inverse, or online optimization is introduced.

\section{Input-Conditioned Detection and Set-Valued Isolation}
\label{sec:detection}

% PLACEHOLDER - Section IV

\subsection{Input-Conditioned Residual Standardization}

The normal distribution of~\eqref{eq:joint_residual} depends on the realized command, command variation, fuzzy operating region, and reference age. Define
\begin{equation}
\begin{aligned}
\bm{\zeta}_k &= [(\bm{\omega}_k^{0})^\mathrm{T} \; \bm{u}_k^\mathrm{T} \; \bm{\upsilon}_k^\mathrm{T} \; \bm{\xi}_k^\mathrm{T} \; \ell_k^0]^\mathrm{T}, \\
\bm{\upsilon}_k &= \bm{u}_k - \bm{u}_{k-1},
\end{aligned}
\label{eq:descriptor}
\end{equation}
where $\bm{\omega}_k^{0}$ contains the firing strengths evaluated along $\bm{z}_k^0$ and its predicted context. For each fuzzy rule, normal training blocks provide a local mean $\bm{\mu}_i(\bm{\zeta}_k)$ and positive elementwise scale $\bm{\sigma}_i(\bm{\zeta}_k)$. Their mixture moments are
\begin{equation}
\begin{aligned}
\bm{\mu}_k &= \sum_i \omega_{i,k}^0 \bm{\mu}_i, \\
\bm{\sigma}_k^2 &= \sum_i \omega_{i,k}^0 (\bm{\sigma}_i^2 + (\bm{\mu}_i - \bm{\mu}_k)^2 + \varepsilon_\sigma \bm{1}).
\end{aligned}
\label{eq:mixture}
\end{equation}
The between-rule term in~\eqref{eq:mixture} prevents uncertainty at a rule boundary from being underestimated. Standardization is elementwise:
\begin{equation}
\widetilde{\bm{e}}_k = (\bm{e}_k - \bm{\mu}_k) \oslash \bm{\sigma}_k.
\label{eq:standardization}
\end{equation}

For each window length $\ell$ in a predeclared finite list, define
\begin{equation}
\varsigma_{k,\ell} = \sum_{j=k-\ell+1}^k \|\widetilde{\bm{e}}_j\|_1.
\label{eq:score}
\end{equation}

% PLACEHOLDER - Section IV-A threshold

A held-out normal calibration split is partitioned by predeclared fuzzy operating cells and reference-age intervals. Within each cell, nonoverlapping blocks calibrate the upper quantile $\tau_{\mathrm{cal}}(\bm{\zeta}_k, \ell)$. The propagated component from~\eqref{eq:beta_recursion}, after the same elementwise scaling, is denoted by $\tau_{\mathrm{prop}}(\beta_{k|s_k^0}, \ell)$. The online threshold is
\begin{equation}
\tau_{k,\ell} = \tau_{\mathrm{cal}}(\bm{\zeta}_k, \ell) + \tau_{\mathrm{prop}}(\beta_{k|s_k^0}, \ell).
\label{eq:threshold}
\end{equation}

Thus, the raw normal region changes with the input, input variation, fuzzy region, and reference age. The calibration term controls the unresolved random remainder, whereas the propagation term records reference drift.

The multiscale statistic is
\begin{equation}
\varsigma_k^{\max} = \max_{\ell \in \{\ell_1, \ldots, \ell_{\max}\}} \frac{\varsigma_{k,\ell}}{\tau_{k,\ell}}.
\label{eq:max_score}
\end{equation}
Its final decision level is calibrated directly on the maximum score, rather than applying separate uncorrected tests to each window. A finite-sample coverage interpretation requires exchangeable calibration and test blocks within each declared cell; under serial dependence, the stated guarantee is blockwise rather than samplewise~\cite{vovk2005,xu2021}.

\subsection{Sufficient Detection Condition}

The following statement clarifies why normal propagation and structured response learning must be used together.

\begin{proposition}[Residual-space detection]
\label{prop:detection}
For a selected window, suppose a normal standardized residual block $\widetilde{\bm{e}}^0$ satisfies $\|\widetilde{\bm{e}}^0\|_1 \leq \tau_{k,\ell}$ on the calibrated normal event. Let a structured perturbation produce the learned response $\Delta\widetilde{\bm{e}}^{g,j,\alpha}$, and suppose the difference between the actual and learned responses is bounded by $\varepsilon_g$. If
\begin{equation}
\|\Delta\widetilde{\bm{e}}^{g,j,\alpha}\|_1 > 2\tau_{k,\ell} + \varepsilon_g,
\label{eq:detection_condition}
\end{equation}
then the corresponding residual block exceeds $\tau_{k,\ell}$ on that event.
\end{proposition}

% PLACEHOLDER - Section IV-B end

The proof is the reverse triangle inequality. If~\eqref{eq:loss_sens_iso} gives $\|\Delta\widetilde{\bm{e}}^{g,j,\alpha}\|_1 \geq \overline{\eta}_g |\alpha|$, a sufficient amplitude condition is
\begin{equation}
|\alpha| > \frac{2\tau_{k,\ell} + \varepsilon_g}{\overline{\eta}_g}.
\label{eq:amplitude}
\end{equation}
Equation~\eqref{eq:amplitude} is a residual-model conclusion, not a physical minimum-fault claim. It shows the two controllable factors: reducing normal propagation lowers $\tau_{k,\ell}$, and increasing the structured response raises $\overline{\eta}_g$.

\subsection{Unique-or-Set-Valued Isolation}

After a detection, the observed residual increment is normalized as in~\eqref{eq:response}. For group $g$, let $\delta_g(k)$ be its smallest $\ell_1$ distance to the response patterns stored from training rollouts. The compatible set is
\begin{equation}
\mathbb{G}_k^{\mathrm{cmp}} = \{g \in \mathbb{G}_b : \delta_g(k) \leq \tau_{\mathrm{iso}}\}.
\label{eq:compatible_set}
\end{equation}
A unique group is returned only when $\mathbb{G}_k^{\mathrm{cmp}}$ contains one element and the distance to the second-best group exceeds a calibrated ambiguity margin.

\begin{proposition}[Sufficient group separation]
\label{prop:separation}
Suppose the observed normalized response lies within $\delta_{\mathrm{obs}}$ of a stored pattern from its generating group. If every pattern from a different group is at an $\ell_1$ distance greater than $2\delta_{\mathrm{obs}}$ from that pattern, nearest-pattern classification returns the generating group.
\end{proposition}

The result follows from the triangle inequality. When its premise is not satisfied,~\eqref{eq:compatible_set} is the appropriate output. The labels $u$, $y$, and $z$ indicate compatibility with an insertion location in the learned model. Without an independently verified mapping, they must not be renamed as unique actuator, sensor, or process faults.

% PLACEHOLDER - Section IV-D

\subsection{Offline and Online Algorithms}

Training, calibration, and testing use disjoint normal trajectory blocks. Algorithm~\ref{alg:offline} summarizes the offline procedure.

\begin{algorithm}
\caption{Normal-Data Offline Design}
\label{alg:offline}
\begin{algorithmic}[1]
\REQUIRE Disjoint normal training and calibration blocks
\STATE Train $\mathcal{H}_\Theta$, $\mathcal{T}_\Theta$, and $\mathcal{D}_\Theta$ by free rollouts using~\eqref{eq:loss_pred}.
\STATE Apply~\eqref{eq:param_metric} and~\eqref{eq:hard_param}; compute~\eqref{eq:bound} and~\eqref{eq:loss_nr}.
\STATE Sample one coordinate from each group in~\eqref{eq:groups}; optimize~\eqref{eq:loss_total} with ordinary forward branches.
\STATE Freeze all model and fuzzy-rule parameters.
\STATE Estimate~\eqref{eq:mixture} and calibrate~\eqref{eq:threshold} on held-out nonoverlapping normal blocks.
\STATE Store response patterns and calibrate the ambiguity margin for~\eqref{eq:compatible_set}.
\end{algorithmic}
\end{algorithm}

Online, the encoder continuously reads the strictly past history, while the normal reference either advances from its accepted anchor or waits for delayed reconfirmation. Algorithm~\ref{alg:online} gives the decision loop.

\begin{algorithm}
\caption{Online Monitoring}
\label{alg:online}
\begin{algorithmic}[1]
\REQUIRE Frozen model, normal statistics, thresholds, and response patterns
\FOR{each sample $k$}
\STATE Compute $\bm{z}_k^{\mathrm{enc}}$ and $\bm{e}_k^y$ from the data branch.
\STATE Propagate the short rollout and the measurement-decoupled reference.
\STATE Form $\bm{e}_k$ by~\eqref{eq:joint_residual} and $\bm{\zeta}_k$ by~\eqref{eq:descriptor}.
\STATE Compute~\eqref{eq:standardization}--\eqref{eq:max_score}.
\IF{$\varsigma_k^{\max}$ exceeds its calibrated decision level}
\STATE Freeze anchor acceptance and report an alarm.
\STATE Return $\mathbb{G}_k^{\mathrm{cmp}}$ by~\eqref{eq:compatible_set}.
\ELSE
\STATE Continue delay-confirmed anchor management.
\ENDIF
\STATE After an alarm, wait $\ell_h$ samples and repeat delay confirmation before accepting a new anchor.
\ENDFOR
\end{algorithmic}
\end{algorithm}

% PLACEHOLDER - Section IV end

The online stage performs encoder and model forward passes, elementwise standardization, vector norms, and table or lightweight-head evaluations. It does not update fuzzy centers, metrics, network weights, or thresholds.

\section{Conclusion}
\label{sec:conclusion}

This paper formulated normal-data-only nonlinear monitoring as three aligned problems: preventing new measurements from updating a normal reference, controlling the resulting free-rollout drift, and making residual-space decisions without physical fault priors. A measurement-decoupled attention--Koopman--T--S reference and a start-indexed residual resolved the training--testing mismatch. The complete fuzzy Jacobian then linked normal propagation to local dynamics, rule overlap, local-model disparity, and scheduling sensitivity through an implementable infinity-norm bound. Structured model-coordinate perturbations complemented the normal propagation objective by increasing residual response and intergroup separation. Finally, input-, fuzzy-region-, and reference-age-conditioned standardization produced a dynamic decision rule with residual-space detection and set-valued isolation conditions. The method does not claim unique physical fault identification: such a claim requires additional actuation measurements, fault-channel knowledge, or a verified mapping from lifted coordinates to physical components.

\appendix

\section{Proof of Theorem~\ref{thm:propagation}}
\label{app:proof1}

From~\eqref{eq:lifted_residual}, add and subtract $\mathcal{T}_\Theta(\bm{z}_k^{\mathrm{enc}}, \bm{u}_k, \bm{\xi}_k; \bm{\chi}_{k|s}^{\mathrm{fr}})$ and then $\mathcal{T}_\Theta(\bm{z}_k^{\mathrm{enc}}, \bm{u}_k, \bm{\xi}_k; \bm{\chi}_k)$. The last difference is $\bm{\nu}_k$, and the context difference is $\bm{\delta}_{k|s}^\chi$. The remaining state difference equals~\eqref{eq:jacobian_integral} $\bm{e}_{k|s}^z$ by the integral mean-value identity. This proves~\eqref{eq:recursion}. Repeated substitution gives~\eqref{eq:propagation}.

\section{Proof of Theorem~\ref{thm:jacobian}}
\label{app:proof2}

Differentiate $\mathcal{T}_\Theta = \sum_i \omega_i \bm{v}_i$. Since $\partial \bm{v}_i / \partial \bm{z} = \bm{A}_i$ and $\partial \bm{\rho} / \partial \bm{z} = \bm{W}_z$,
\begin{equation}
\bm{\Gamma}_{z,k} = \sum_i \omega_{i,k} \bm{A}_i + \sum_i \bm{v}_{i,k} (\nabla_{\bm{\rho}} \omega_{i,k})^\mathrm{T} \bm{W}_z.
\end{equation}
Substitution of~\eqref{eq:gradient}, together with $\sum_i \omega_i (\nabla \varphi_i - \sum_j \omega_j \nabla \varphi_j) = 0$, permits replacing $\bm{v}_{i,k}$ by $\bm{v}_{i,k} - \overline{\bm{v}}_k$, which yields~\eqref{eq:complete_jacobian}. For vectors $\bm{a}$ and $\bm{b}$, $\|\bm{a}\bm{b}^\mathrm{T}\|_\infty = \|\bm{a}\|_\infty \|\bm{b}\|_1$. The triangle inequality and submultiplicativity then give~\eqref{eq:bound}.


\bibliographystyle{IEEEtran}
\bibliography{refs}

\end{document}
